\subsection{I-35. A smaller contraction invariant, its exact composition rule, and the rank-one lift obstruction}
\label{I-35}

\paragraph{Scope.}
This result concerns the rank-$r$ outer matrix products left after
factoring higher-rank resets. It separates three assertions:
preserving all constant scalar boundary contractions of one bridge,
composing bridge substitutions with independent parameters, and
synchronizing those parameters. The first is strictly weaker than
preserving coefficient span or generic matrix rank. The second has
an exact positive proof. The third does not follow, and explicit
counterexamples below identify the missing quantifier. No general
ROABP hitting set is claimed.

\paragraph{1. The exact rank-one-annihilator criterion.}
Let $M(\mathbf X)$ be an $r\times r$ polynomial matrix over a field
$K$, and let $U\subseteq\operatorname{Mat}_r(K)$ be the span of its
coefficient matrices. After a polynomial substitution $G$, write
$U_G$ for the coefficient-matrix span of $M(G)$; then $U_G\subseteq U$.
All statements can instead be made over a specified working extension.
For rational substitutions, clear a common nonzero polynomial
denominator first. The span is independent of this clearing choice:
a constant functional annihilates the cleared vector if and only if
it annihilates the original rational vector.

Use the pairing $\langle C,A\rangle=\operatorname{tr}(C^{\mathsf T}A)$,
and let $\Sigma_1$ be the set of matrices of rank at most one. Then
\[
 \begin{gathered}
 uM(\mathbf X)v\ne0\ \Longrightarrow\ uM(G)v\ne0
       \quad\hbox{for every constant row }u\hbox{ and column }v
 \\
 \Longleftrightarrow\qquad
 U_G^\perp\cap\Sigma_1=U^\perp\cap\Sigma_1.
 \end{gathered}
 \tag{I.75}
\]
Indeed, $uAv=\langle u^{\mathsf T}v^{\mathsf T},A\rangle$, and every
nonzero rank-one matrix has this form. A scalar contraction is the
zero polynomial exactly when its functional annihilates the entire
coefficient space. Since $U_G\subseteq U$, one inclusion of the two
intersections is automatic and the other is precisely nonvanishing
preservation. For preservation after every scalar extension, impose
this equality over an algebraic closure. Thus the target is equality
of base loci on the Segre variety of boundary pairs, rather than
equality of the entire linear spaces or matrix image cones.

\paragraph{2. An explicit optimal test space of dimension $2r-1$.}
Index rows and columns by $0,\ldots,r-1$. Let $H_k$ have entry one
when $i+j=k$ and zero otherwise, for $0\le k\le2r-2$. Their disjoint
nonempty supports make them linearly independent. Put
\[
 a(T)=(1,T,\ldots,T^{r-1})^{\mathsf T},\qquad
 H(T)=a(T)a(T)^{\mathsf T}=\sum_{k=0}^{2r-2}T^kH_k.
 \tag{I.76}
\]
For any nonzero constant $u,v$,
\[
 uH(T)v=
   \left(\sum_{i=0}^{r-1}u_iT^i\right)
   \left(\sum_{j=0}^{r-1}v_jT^j\right)\ne0.
 \tag{I.77}
\]
This is a product of two nonzero polynomials over a domain, in every
characteristic. Consequently the $(2r-1)$-dimensional Hankel space
$\mathcal H_r=\operatorname{span}\{H_k\}$ has no nonzero rank-one
annihilator. It satisfies (I.75) with $U=\operatorname{Mat}_r(K)$,
although $\mathcal H_r$ is a proper subspace when $r\ge2$.
Evaluating $H(T)$ at $2r-1$ distinct field values spans this space by
Vandermonde interpolation; an extension may supply the values.

The dimension $2r-1$ is optimal over an algebraically closed field.
This is known completely-entangled-subspace background, not a new
claim: Parthasarathy (2004), Propositions~1.3--1.4 and Theorem~1.5
(P-B16), gives this two-factor construction and optimal codimension
over $\mathbb C$. The algebraic-closure formulation follows from
the same Segre dimension argument; its checked proof is retained
in the separate background audit. The lower bound need not hold
when only rational boundary pairs over a nonclosed field are tested.

This is an actual matrix-bridge substitution. Take
$M(\mathbf X)=(X_{ij})$, whose coefficient span is all of
$\operatorname{Mat}_r(K)$, and substitute $X_{ij}=T^{i+j}$.
The generic matrix rank drops from $r$ to one, while (I.77) preserves
every nonzero constant boundary contraction. Moreover this generic
matrix is realized by an invertible width-$2r$ interior word:
\[
 \prod_{i,j}\bigl(I_{2r}+X_{ij}E_{i,r+j}\bigr)
     =\begin{pmatrix}I_r&M(\mathbf X)\\0&I_r\end{pmatrix}.
 \tag{I.78}
\]
Here one may relabel matrix-unit indices from one instead of zero.
All products of two off-diagonal generators vanish, so the identity
holds in every order. Selecting the upper-right block uses constant
full-rank rectangular endpoints. Thus generic rank preservation is
strictly stronger than scalar boundary preservation even for an
actual bridge with invertible interior factors.

More generally, after scalar extension to an algebraic closure,
every coefficient space $U$ has a subspace $W\subseteq U$ of
dimension at most $2r-1$ with the same rank-one annihilators.
To prove this, remove the base locus from the Segre
variety and map a boundary pair to its nonzero evaluation functional
in $\mathbb P(U^*)$. The closure $Y$ of this image has dimension at
most $2r-2$. If $\dim U\le2r-1$, take $W=U$. Otherwise set
$D=\dim U$ and $c=2r-1$. The Grassmannian of $(D-c)$-dimensional
vector subspaces of $U^*$ has dimension $c(D-c)$. Requiring such
a subspace to contain one fixed projective point imposes $c$
conditions: its fiber is a Grassmannian of dimension
$c(D-c-1)$. The incidence variety of subspaces meeting $Y$
therefore has dimension at most
$\dim Y+c(D-c-1)<c(D-c)$. Its image is closed because $Y$ is
projective, so a nonempty open set of subspaces misses $Y$.
The annihilator of any one of them defines the desired $W$.
This proves existence of a small contraction test space, not
an explicit coefficient-oblivious substitution realizing it for
an arbitrary bridge. The Hankel construction realizes it for the
full generic matrix example only. No field-of-definition or
effective-construction bound is asserted for the general existence
statement.

\paragraph{3. The valid composition statement and its independent-seed cost.}
Let $B_j(\mathbf X_j)$, $1\le j\le s$, be square matrices with
pairwise disjoint variable sets. Suppose each substitution
$\mathbf X_j=G_j(\mathbf S_j)$ preserves every constant scalar
boundary contraction of $B_j$, and does so after every field extension.
If the seed sets $\mathbf S_j$ are algebraically independent of one
another and of all remaining variables, then
\[
 uB_1(\mathbf X_1)\cdots B_s(\mathbf X_s)v\ne0
 \ \Longrightarrow\
 uB_1(G_1)\cdots B_s(G_s)v\ne0.
 \tag{I.79}
\]
Replace the matrices successively. At step $j$, take as coefficient
field the rational function field in all unprocessed variables and
all previously introduced seeds. The product on the left of $B_j$
is a constant row over this field, and the product on its right is
a constant column. The current scalar polynomial is nonzero, so
the hypothesis over that extension field preserves its nonzeroness
when the fresh seed set $\mathbf S_j$ is introduced. Induction proves
(I.79). This argument requires no equality of coefficient spaces.

If adjacent bridges share cut variables $\mathbf Z$, first work
over the field $K(\mathbf Z)$ and leave those variables unassigned.
Assume the bridge preservation hypothesis over this field and all
its extensions. The same proof then applies to substitutions of
their disjoint interior variables. A later specialization of $\mathbf Z$ must still preserve
the resulting nonzero numerator and any relevant denominators.
Nothing in (I.79) supplies that specialization automatically.

For a prescribed decomposition into $s$ blocks, suppose each map
is polynomial over $K$ in its seeds and the retained cut variables,
uses at most $d$ seeds, and has total coordinate degree at most $L$.
If $z$ shared cut variables remain free, then the resulting
original scalar polynomial of total degree at most $D$ has at most
$sd+z$ parameters and total degree at most $D\max\{1,L\}$.
A full grid therefore costs at most
\[
 \bigl(D\max\{1,L\}+1\bigr)^{sd+z}.
 \tag{I.80}
\]
This displays the cost rather than concealing it in a genericity
assertion. The argument presupposes the decomposition; it does not
produce a uniform unknown-order sampler or identify the reset
locations from a scalar black box. For maps with arbitrary
$K(\mathbf Z)$ coefficients instead, numerator and denominator
degrees require separate bounds; (I.80) is not asserted for them.

\paragraph{4. Why constant contraction preservation does not synchronize.}
Even preserving each entire local coefficient space does not justify
using the same seed in different blocks. For example, with
$u=e_1^{\mathsf T}$ and $v=e_2$,
\[
 u(I+XE_{12})(I-YE_{12})v=X-Y\ne0.
 \tag{I.81}
\]
Each substitution $X=T$ or $Y=T$ separately preserves the full
coefficient space of its invertible matrix factor. Simultaneously
making both substitutions kills the scalar contraction.

The stronger Hankel example makes the boundary quantifier explicit.
For $r=2$, the map $H(T)$ in (I.76) preserves every nonzero constant
boundary contraction, but
\[
 (-T,1)H(T)e_1=0.
 \tag{I.82}
\]
With an independent boundary parameter $S$, the corresponding
expression is $T-S\ne0$. Diagonal identification $S=T$ creates the
cancellation. Treating all constant boundary pairs over extensions
is not a license to identify a new seed with a preexisting element
of the coefficient field.

In fact an unrestricted same-seed rational-boundary condition would
be impossible for $r\ge2$. Given any matrix $A(\mathbf S)$ over
$K(\mathbf S)$, there are nonzero $u(\mathbf S),v(\mathbf S)$ with
$uAv=0$: choose a nonzero row and then a nonzero vector in the
kernel of its resulting row functional. For the generic matrix
$M(\mathbf X)$, however, $uM(\mathbf X)v$ is nonzero over
$K(\mathbf S)$ because at least one product $u_iv_j$ is nonzero.
Thus a viable synchronization theorem must use restrictions on the
boundary functions actually induced by neighboring bridges. Asking
for all rational functions of the same seed is too strong.

\paragraph{5. A rank-one multiplicative lift cannot retain an entry for $r\ge2$.}
There is also a precise obstruction to turning the higher-rank
outer product into scalar product closure by a multiplicative lift.
Let $r\ge2$, let $w\ge r$, and embed $r\times r$ matrices as the
upper-left corner $\iota(B)$ of $w\times w$ matrices. Put
$P=\iota(I_r)$. Suppose a multiplicative map
\[
 \rho:\operatorname{Mat}_w(K)\longrightarrow\operatorname{Mat}_N(E)
\]
is given over some extension $E$, and suppose $\rho(P)$ has rank
one. There cannot be a fixed $E$-linear readout $\Lambda$ satisfying
\[
 \Lambda(\rho(\iota(B)))=B_{11}
 \quad\hbox{for every }B\in\operatorname{GL}_r(K).
 \tag{I.83}
\]
No polynomiality or linearity of $\rho$ is needed for this obstruction;
multiplicativity and the fixed linear readout suffice.

Indeed $Q=\rho(P)$ is a rank-one idempotent. Since
$P\iota(B)P=\iota(B)$, every lifted corner matrix lies in
$Q\operatorname{Mat}_N(E)Q=EQ$ and has the form $\sigma(B)Q$.
The identity case of (I.83) gives $\Lambda(Q)=1$, so
$\sigma(B)=B_{11}$. Multiplicativity would then make $B\mapsto B_{11}$
a multiplicative function on $\operatorname{GL}_r(K)$. But for
\[
 B=I_r+E_{12},\qquad C=I_r+E_{21},
\]
one has $B_{11}=C_{11}=1$ and $(BC)_{11}=2\ne1$, in every
characteristic. This is a contradiction. Multiplying the proposed
readout by any integral power of $\det B$ does not evade it, because
$B,C,BC$ all have determinant one.

Exterior powers are consistent with this obstruction: at a rank-$r$
reset, the $r$th exterior power has rank one and retains maximal-minor
data, which acts by the determinant on an invertible $r$-dimensional
corner. It cannot also recover an arbitrary fixed entry there by
such a readout. This excludes the stated functorial scalarization,
not all nonfunctorial or coefficient-dependent constructions.

\paragraph{Remaining concrete gate.}
The relevant smaller target is the Segre-annihilator condition for
the actual contractions, together with a synchronization theorem
for the restricted boundary functions produced by adjacent reset
bridges. The independent-seed statement and its cost are proved
above. The Hankel bridge proves that this target can preserve scalar
information while destroying generic matrix rank. Neither fact
establishes an efficient synchronized map for general rank-$r$
outer words. That last assertion remains open in this report.
\hfill$\square$

\paragraph{Background and provenance.}
The optimal Hankel test-space statement is the known two-factor
case of Parthasarathy (2004), Propositions~1.3--1.4 and Theorem~1.5,
preprint pp.~4--6, already recorded in P-B16. The complete local
proofs were independently reopened and read for this audit.
The projective facts used here are the Segre embedding, Stacks
Project Authors (2026), Lemma~27.13.6, tag~01WD, and the projective
intersection dimension theorem, Lemma~33.34.3, tag~0B2R. Their
primary statements and displayed proofs were read. The Hankel
identity, actual unipotent bridge realization, composition rule,
synchronization counterexamples, and multiplicative-corner
obstruction are proved directly above. No historical priority
claim or general derandomization consequence is made.
